%%%%

%%%   Самарский государственный технический университет

%%%   Кафедра прикладной математики и информатики

%%%

%%%   Преамбула для статьи в журнал "Вестник Самарского государственного

%%%   технического университета. Серия: «Физико-математические науки»"

%%%   Ориентирована под MikTEx 2.4 for Win32

%%%

%%%   Хотя сам журнал собирается под GNU/Linux на дистрибутиве TeXLive



\documentclass[11pt,a4paper,twoside]{article}



\usepackage[cp1251]{inputenc}

\usepackage[T2A]{fontenc}

\usepackage[english,russian]{babel}

\usepackage{samgtubib}

\usepackage[dvips]{graphicx}

\usepackage[multiple]{footmisc}



\usepackage{samgtu}

\renewcommand{\VestnikYear}{2009} % Год издания Вестника

\renewcommand{\VestnikNumber}{2\,(19)} % Номер Вестника

\renewcommand{\VestnikMonth}{Ноябрь} % Месяц выхода Вестника

\renewcommand{\VestnikData}{01.11.09} % Дата подписания Вестника в печать

\renewcommand{\VestnikUPL}{26,64} % Усл. печ. л.

\renewcommand{\VestnikUIL}{25,73} % Уч.-изд. л.

\renewcommand{\VestnikRegNumber}{479/09} % Регистрационный номер

\renewcommand{\VestnikOrder}{994} % Номер заказа







%

%\usepackage{textcomp}

%

%\begin{document}

%





\renewcommand{\firstpage}{291}

\renewcommand{\lastpage}{296}



%\ShortPartition{Математическое моделирование}{ }

%

\udk{517.958+519.633 } \msc{80A17}





\title{Дополнительные граничные условия в задачах теплопроводности с переменными физическими свойствами среды}% Полный заголовок

{Дополнительные граничные условия в задачах теплопроводности \dots}% Заголовок, который идёт в верхний колонтитул



\author{Е.\,В.~Ларгина, Е.\,В.~Стефанюк, В.\,В.~Котов}{Л\,а\,р\,г\,и\,н\,а~Е.\,В., С\,т\,е\,ф\,а\,н\,ю\,к~Е.\,В., К\,о\,т\,о\,в~В.\,В.}



\institute{\samgtu.}



\email{E-mails}{larginaevgenya@mail.ru, stef-kate@yandex.ru, cx65@list.ru}



\otherinf{\fullauthor{Евгения Валериевна Ларгина} \position {аспирант} \dept {каф. теоретических основ теплотехники и гидромеханики}

\fullauthor{Екатерина Васильевна Стефанюк} \regalia{к.т.н., доцент} \position{доцент} \dept {каф. теоретических основ теплотехники и гидромеханики}

\fullauthor{Вячеслав Васильевич Котов} \position{инженер} \dept {каф. теоретических основ теплотехники и гидромеханики}

}











\keywords{приближённое аналитическое решение, дополнительные граничные условия, собственные числа, уравнение Штурма"--~Лиувилля}



\workabstract{Разрабатывается направление получения приближенных аналитических решений краевых задач математической физики,

основанное на использовании дополнительных граничных условий, позволяющих в модельном (апроксимационном) представлении

аналитического решения определять любое число его слагаемых и таким образом получать решения с заданной степенью точности.

Физический смысл дополнительных граничных условий состоит в выполнении основного дифференциального уравнения и производных от

него различного порядка в граничных точках области, что приводит к выполнению уравнения и внутри области с точностью, зависящей

от числа приближений.}



\engtitle{Additional Boundary Condition in Heat Conductivity Problems with Variables Medium Physical Properties}



\engauthor{E.\,V.~Largina, E.\,V.~Stefanjuk, V.\,V.~Kotov}{L\,a\,r\,g\,i\,n\,a~E.\,V., S\,t\,e\,f\,a\,n\,j\,u\,k~E.\,V.,

K\,o\,t\,o\,v~V.\,V.}



\otherenginfm{\fullauthor{Evgeniya V. Largina} \position{Postgraduate Student} \dept{Dept. of Theoretical Basis of Heat Engineering and Flow Mechanics} 

\fullauthor{Ekaterina V. Stefanyuk}  \regalia{Ph.\,D. (Tech.)} \position{Associate Professor}  \dept{Dept. of Theoretical Basis of Heat Engineering and Flow Mechanics} 

\fullauthor{Vyacheslav V Kotov} \position{Engineer} \dept{Dept. of Theoretical Basis of Heat Engineering and Flow Mechanics} 

}





\enginstitute{\engsamgtu.}



\engabstract{The field of obtaining approximated analytical solutions of mathematical physics boundary problems is developed. It is based on additional boundary conditions, allowing to define any item number of model (approximated) analytical solutions and to obtain solutions with the required accuracy degree. Physical meaning of additional boundary conditions lies in implementation of the basic differential equation and its different order derivatives at boundary points of the field. It leads to implementation of differential equation within the given field with number of approximations dependant on the accuracy.}



\engkeywords{approximated analytical solutions, additional boundary conditions, eigenvalue, Sturm--Liouville equation}



\date{28/VI/2010}

\revisiondate{15/IX/2010}





\maketitle  \small



Точные аналитические решения нестационарных задач теплопроводности с переменными по пространственным координатам физическими

свойствами среды в настоящее время получены лишь для одномерного полупространства \cite{larg10:1}. При этом решения, как правило,

выражаются сложными функциональными зависимостями, использование которых в инженерной практике оказывается весьма

затруднительным. В связи с этим актуальной является проблема получения хотя бы приближенных аналитических решений с точностью,

достаточной для инженерных приложений. Для нахождения таких решений весьма эффективным является метод, основанный на

использовании дополнительных граничных условий \cite{larg10:2,larg10:3} и позволяющий в ряде случаев получать решения, практически эквивалентные

точным.



В качестве конкретного примера использования этого метода найдем решение

нестационарной задачи теплопроводности для бесконечной пластины при

экспоненциальном изменении коэффициента теплопроводности от пространственной

координаты

\begin{equation}

\label{larg10:eq1}

\lambda (x)=\lambda_0 \exp (-mx),

\end{equation}

$\lambda_0$ "--- величина коэффициента теплопроводности при $x=0$;

$m>0$ "--- коэффициент, характеризующий интенсивность изменения $\lambda (x)$.



Из соотношения (\ref{larg10:eq1}) следует, что при фиксированном значении величины $m$

теплопроводность пластины с увеличением координаты уменьшается.



Математическая постановка задачи для симметричных граничных условий первого

рода в данном случае имеет вид:

\begin{gather}

\label{larg10:eq2}

c\rho \frac{\partial T(x,\tau)}{\partial \tau }=\frac{\partial }{\partial

x}\Bigl[ \lambda (x)\frac{\partial T(x,\tau )}{\partial x} \Bigr] \quad

(\tau >0, \quad 0\le x\le \delta );\\

\label{larg10:eq3}

T(x,\,0)=T_0;\\

\label{larg10:eq4}

T(\delta,\,\tau )=T_\text{ст};\\

\label{larg10:eq5}

\frac{\partial T(0,\,\tau )}{\partial x}=0,

\end{gather}

где $c$ "--- теплоемкость, $\rho $ "--- плотность, $T$ "--- температура, $x$ "---

координата, $\tau $ "--- время, $T_0 $ "--- начальная температура,

$T_\text{ст}$ "--- температура стенки при $x=\delta $.



Введём следующие безразмерные переменные:

\begin{equation}

\label{larg10:eq6}

\Theta =({T-T_\text{ст}})/({T_0 -T_\text{ст}},) \quad

\xi ={ x}/{\delta}, \quad

{\tt Fo}={a \tau}/{\delta^2}.

\end{equation}

Задача (\ref{larg10:eq2})--(\ref{larg10:eq5}) с учётом (\ref{larg10:eq6}) примет вид

\begin{gather}

\label{larg10:eq7}

\frac{\partial \Theta (\xi,\,{\tt Fo})}{\partial {\tt Fo}}=\frac{\partial }{\partial \xi

}\left[ e^{-\nu \xi }\frac{\partial \Theta (\xi,\,{\tt Fo})}{\partial \xi}\right]

\quad

({\tt Fo}>0; \, 0 \le x\le 1); \\

\label{larg10:eq8}

\Theta (\xi,\,0)=1;\\

\label{larg10:eq9}

\Theta (1,\,{\tt Fo})=0;\\

\label{larg10:eq10}

\frac{\partial \Theta (0,\,{\tt Fo})}{\partial \xi}=0,

\end{gather}

где $\nu = m \delta$.

Решение задачи (\ref{larg10:eq7})--(\ref{larg10:eq10}) разыскивается в виде произведения двух функций:

\begin{equation}

\label{larg10:eq11}

\Theta (\xi,\,{\tt Fo})=\varphi ({\tt Fo})\Psi (\xi ).

\end{equation}

Подставляя (\ref{larg10:eq11}) в (\ref{larg10:eq7}), находим

\begin{gather}

\label{larg10:eq12}

\frac{d\varphi ({\tt Fo})}{d{\tt Fo}} + \mu \varphi ({\tt Fo})=0;\\

\label{larg10:eq13}

\frac{d}{d\xi}\Bigl[ e^{-\nu \xi }\frac{d \Psi (\xi )}{d\xi} \Bigr]+\mu \Psi (\xi )=0,

\end{gather}

где $\mu $ "--- некоторая постоянная.

Решение уравнения (\ref{larg10:eq12}) известно и имеет вид

\begin{equation}

\label{larg10:eq14}

\varphi ({\tt Fo})=A \exp (-\mu {\tt Fo}),

\end{equation}

где $A$ "--- неизвестный коэффициент.



Найдём приближенное аналитическое решение уравнения Штурма"--~Лиувилля (\ref{larg10:eq13}), для которого

граничные условия согласно (\ref{larg10:eq9}), (\ref{larg10:eq10}) будут такими:

\begin{gather}

\label{larg10:eq15}

\Psi (1)=0;\\

\label{larg10:eq16}

\Psi'(0)=0.

\end{gather}

Решение задачи (\ref{larg10:eq13}), (\ref{larg10:eq15}), (\ref{larg10:eq16}) ищется в виде ряда

\begin{equation}

\label{larg10:eq17}

\Psi (\mu ,\xi )=\sum\limits_{i=1}^n a_i N_i (\xi ),

\end{equation}

где $a_i$ "--- неизвестные коэффициенты, $N_i (\xi)=\cos (r\pi {\xi}/{2})$  "--- координатные функции, $r=2i-1$.

Благодаря принятой системе координатных функций $N_i (\xi)$ соотношение (\ref{larg10:eq17}) точно удовлетворяет граничным условиям (\ref{larg10:eq15}),

(\ref{larg10:eq16}).



Для нахождения решения в первом приближении исходя из (\ref{larg10:eq16}) введём

следующее дополнительное граничное условие:

\begin{equation}

\label{larg10:eq18}

\Psi (0)=1.

\end{equation}

Ограничиваясь одним членом ряда в соотношении (\ref{larg10:eq17}), для нахождения

коэффициента $a_1$ подставим (\ref{larg10:eq17}) в (\ref{larg10:eq18}):

$[a_1 \cos ({\pi }\xi/2 )]_{\xi =0} =1$, откуда $a_1 = 1$.

Соотношение (\ref{larg10:eq17}) в первом приближении принимает вид

\begin{equation}

\label{larg10:eq19}

\Psi (\xi )=\cos ({\pi }\xi/2).

\end{equation}



Для определения первого собственного числа составим невязку уравнения (\ref{larg10:eq13}) и

проинтегрируем её в пределах от $0$ до $1$:

\begin{equation}

\label{larg10:eq20}

\int_0^1 \left( -\frac{\pi }{2e^\xi }\sin \frac{\pi \xi

}{2}-\frac{\pi ^2}{4e^\xi }\cos ({\pi \xi }/{2})+\mu \cos ({\pi \xi }/{2}) \right) d\xi =0.

\end{equation}

Вычисляя в (\ref{larg10:eq20}) интегралы, находим $\mu _1 =0,908$.



Для получения решения во втором приближении введём ещё одно дополнительное

граничное условие, получаемое из уравнения (\ref{larg10:eq13}) при $\xi =0$:

\begin{equation}

\label{larg10:eq21}

\Psi''(0)=-\mu .

\end{equation}

Подставим (\ref{larg10:eq17}), ограничиваясь двумя членами ряда, в граничные условия (\ref{larg10:eq18}),

(\ref{larg10:eq21}) для определения неизвестных коэффициентов $a_1 $ и $a_2 $ и получим

систему двух алгебраических линейных уравнений, которая имеет следующее решение:

\begin{equation}

\label{larg10:eq22}

a_1 = 9 \pi^2 - \frac{4\mu}{8 \pi ^2}, \quad

a_2 = 4 \mu - \frac{\pi ^2}{8 \pi^2}.

\end{equation}

Подставляя (\ref{larg10:eq22}) в (\ref{larg10:eq17}), находим

\begin{equation}

\label{larg10:eq23}

\Psi(\xi )=\left[ 9 \pi^2 - \frac{4 \mu}{8\pi^2} \right] \cos \left({\pi \xi }/{2}\right) + 4 \mu - \frac{\pi ^2}{8\pi ^2}

\cos \left( {3\pi \xi }/{2}\right).

\end{equation}

Подставляя (\ref{larg10:eq23}) в (\ref{larg10:eq13}) и вычисляя интегралы в пределах от $0$ до $1$, относительно первых двух собственных чисел получаем следующее характеристическое уравнение:

\begin{equation}

\label{larg10:eq24}

-0{,}043067512 \mu^2 + 0{,}850297239 \mu -0{,}86664=0,

\end{equation}

корни которого представляют собой два первых собственных числа

уравнения Штурма"--~Лиувилля (\ref{larg10:eq13}): $\mu _1 \approx 1,078$, $\mu_2 \approx 18,665$.





Подставляя (\ref{larg10:eq14}), (\ref{larg10:eq17}) в (\ref{larg10:eq11}), для каждого собственного числа будем иметь

частные решения вида

\begin{equation}

\label{larg10:eq27}

\Theta_i (\xi, {\tt Fo})=A_i \Psi_i (\mu_i,\,\xi)\exp (-\mu_i {\tt Fo}).

\end{equation}



Каждое частное решение вида (\ref{larg10:eq27}) точно удовлетворяет основным граничным условиям (\ref{larg10:eq15}), (\ref{larg10:eq16}) и~приближённо (в~зависимости от числа членов ряда (\ref{larg10:eq17})\,) удовлетворяет уравнению (\ref{larg10:eq7}). 

Однако ни одно из них, в том числе и их сумма

\begin{equation}

\label{larg10:eq28}

\Theta_n (\xi ,\,{\tt Fo})=\sum \limits_{i=1}^n A_i \Psi_i (\mu_i,\,\xi ) \exp (-\mu_i {\tt Fo}),

\end{equation}

не удовлетворяют начальному условию (\ref{larg10:eq8}). Для того чтобы соотношение (\ref{larg10:eq28})

удовлетворяло начальному условию (\ref{larg10:eq8}), составляется его невязка и требуется

ортогональность невязки ко всем координатным функциям, т.\,е.

\begin{equation}

\label{larg10:eq291}

\int _0^1 \left[ \sum \limits_{i=1}^n A_i \Psi _i (\mu_i, \xi)-1 \right] N_j (\xi ) d\xi =0, \quad j=1, 2, \dots, n.

\end{equation}



Вычисляя интегралы в \eqref{larg10:eq291}, для нахождения коэффициентов $A_i$ получим

систему из $n$ алгебраических линейных уравнений по количеству собственных чисел.



Соотношение \eqref{larg10:eq291} для второго приближения запишется так:

$$

\int_0^1 \left[ A_1 \Psi _1 (\mu_1,\,\xi ) + A_2 \Psi _2 (\mu_2

,\xi)-1 \right] N_1 (\xi )d\xi =0, 

$$

$$

\int_0^1 \left[ A_1 \Psi _1 (\mu _1,\,\xi ) + A_2 \Psi_2 (\mu_2

,\xi)-1 \right] N_2 (\xi )d\xi =0.

$$

Система подставлением собственных функций, соответствующих собственным числам,  сводится к системе двух алгебраических линейных уравнений, решение которой $A_1 =1,251$, $A_2 =- 0,402$.





Для получения решения в шестом приближении необходимо  к~граничным условиям (\ref{larg10:eq15}), (\ref{larg10:eq16}), (\ref{larg10:eq18}), (\ref{larg10:eq21}) добавить следующие дополнительные граничные условия:

\begin{equation}

\label{larg10:eq30}

\begin{array}{l}

\Psi '''(0)=0, \qquad \Psi^{\rm IV}(0)=\mu^2 - 3\mu \nu^2, \qquad \Psi^{\rm V}(0)=0, \\

\Psi^{\rm VI}(0)=-\mu^3 + 23\mu^2 \nu^2 - 5 \mu \nu^4; \qquad \Psi^{\rm VII}(0)=0, \\

\Psi^{\rm VIII}(0)=\mu^4 - 86 \mu^3 \nu^2 + 201 \mu^2 \nu^4 - 7\mu \nu^6, \quad \Psi^{\rm IX}(0)=0,\\

\Psi^{\rm X} (0) = -\mu^5 + 230 \mu^4 \nu^2 + 2307 \mu^3 \nu^4 + 1291 \mu^2 \nu^6 - 9\mu \nu^8.

\end{array}

\end{equation}



Дополнительные граничные условия (\ref{larg10:eq30}) так же, как и (\ref{larg10:eq21}),  получаются из дифференциального уравнения (\ref{larg10:eq13}) путём его дифференцирования применительно к точке $\xi =0$.

Отметим, что соотношение (\ref{larg10:eq17}) удовлетворяет дополнительным однородным граничным

условиям (\ref{larg10:eq30}) %, полученным из уравнения (\ref{larg10:eq13}) 

с нечётными производными,  а также условиям (\ref{larg10:eq15}), (\ref{larg10:eq16}).



%Для определения коэффициентов $a_i $ $(i= \overline {1,\,6} )$ используются

%пять дополнительных граничных условий, полученных из дифференциального

%уравнения (\ref{larg10:eq13}) с чётными производными, т.\,е. используются условия:

%\begin{gather*}

%\Psi^{II}(0)=-\mu; \, \Psi^{IV}(0)=\mu^2 - 3\mu \nu^2; \, \Psi^{VI}(0)=-\mu^3 + 23 \mu^2 \nu^2 - 5\mu \nu^4;\\

%\Psi^{VIII}(0) = \mu^4 - 86 \mu^3 \nu^2 + 201\mu^2 \nu^4 - 7\mu \nu^6;\\

%\Psi^X(0)=-\mu^5 + 230 \mu^4 \nu^2 + 2307 \mu^3\nu^4 + 1291 \mu^2 \nu^6 - 9 \mu \nu^8\,

%\end{gather*}

%а также условие (\ref{larg10:eq18}).

Подстановка (\ref{larg10:eq17}) в граничные условия (\ref{larg10:eq30}) с~чётными производными  и в \eqref{larg10:eq18}  даёт систему шести алгебраических линейных уравнений, разрешая которую относительно $a_i$, находим

$$

\begin{array}{l}

a_1 {=} {-}1,234 {\cdot }10^{-8} \mu^5 + 1,152 {\cdot} 10^{-5} \mu^4 - 2,914 {\cdot}

10^{-3} \mu^3 + 0,274 {\cdot} 10^{-2} \mu^2 - 0,101 \mu + 1,221;\\

a_2 {=} 2,644 {\cdot} 10^{-8}\mu^5 - 2,41 {\cdot} 10^{-5}\mu ^4{+}5,84{\cdot} 10^{-3}\mu

^3{-}0,498{\cdot} 10^{-2}\mu ^2{-}0,139\mu {-}0,29;\\

a_3 {=} {-}2,204\cdot 10^{-8}\mu ^5{+}1,926{\cdot} 10^{-5}\mu ^4{-}4,254{\cdot}

10^{-3}\mu ^3{+}0,296{\cdot} 10^{-2}\mu ^2{-}0,047\mu {+}0,087;\\

a_4{=}1,028{\cdot}10^{-8}\mu ^5{-}8,382{\cdot} 10^{-6}\mu ^4{+}1,614{\cdot} 10^{-3}\mu

^3{-}8,707{\cdot} 10^{-2}\mu ^2{+}0,011\mu {+}0,02;\\

a_5 {=} {-}2,644\cdot 10^{-9}\mu ^5{+}1,194{\cdot} 10^{-6}\mu ^4{-}3,161{\cdot}

10^{-4}\mu ^3{+}1,478{\cdot} 10^{-2}\mu ^2{-}0,0018\mu{ +}0,003;\\

a_6 {=}2,93{\cdot} 10^{-10}\mu ^5{-}1,87{\cdot}10^{-7}\mu ^4{+}0,266{\cdot} 10^{-6}\mu

^3{+}1,149{\cdot} 10^{-3}\mu ^2{+}0,138{\cdot}10^{-3}\mu {-}0,002.

\end{array}

$$



После подстановки найденных значений $a_i$ в (\ref{larg10:eq17}) составляется интеграл взвешенной невязки уравнения (\ref{larg10:eq13}):

\begin{equation}

\label{larg10:eq31}

\int _0^1 

\left[

\frac{d^2}{d\xi^2} \Bigl( \sum \limits_{i=1}^6 a_i \cos (r {\pi }\xi/{2}) \Bigr) + \mu \sum\limits_{i=1}^6 a_i \cos (r {\pi }\xi/2) \right]

d\xi =0,\quad r=2i-1.

\end{equation}

Вычисляя  интегралы в (\ref{larg10:eq31}), относительно собственных чисел $\mu $ получаем

алгебраическое уравнение шестой степени, корни которого следующие:

$\mu_1 = 1,314$,

$\mu_2 = 16,734$,

$\mu_3 = 45,061$,

$\mu_4 = 84,603$,

$\mu_5 = 181,998$,

$\mu_6 = 588,989$.





Соответствующие каждому собственному числу собственные функции находятся из

(\ref{larg10:eq17}). Подставляя собственные функции в (\ref{larg10:eq11}), получаем

\begin{equation}

\label{larg10:eq331}

\Theta (\xi,\,{\tt Fo}) = \sum\limits_{i=1}^6 A_i a_i (\mu_i)\cos (r{\pi }\xi/2 )\cdot  \exp (-\mu_i {\tt Fo}),\quad r=2i-1.

\end{equation}



Для нахождения неизвестных коэффициентов $A_i$ составляется невязка начального условия (\ref{larg10:eq8})

и требуется ортогональность невязки к каждой координатной функции, т.\,е.

\begin{equation}

\label{larg10:eq341}

\int _0^1 \biggl[ \sum\limits_{i=1}^6 A_i a_i (\mu_i) \cos (r {\pi} \xi/2) - 1 \biggr] \cos (j{\pi }\xi/2) d \xi =0,

\quad j=r= 2i-1.

\end{equation}

%

%Для данного конкретного случая при таких простых начальных условиях (не

%зависящих от координаты $\xi$) процесс определения коэффициентов $A_i$

%может быть существенно упрощён (по сравнению со случаем их нахождения для

%двух приближений, см.~систему уравнений (\ref{larg10:eq29})).



Ввиду ортогональности косинусов система уравнений \eqref{larg10:eq341} приводится к системе

уравнений, в которой неизвестные разделяются (в каждое уравнение входит лишь

одна неизвестная величина $A_i$). В связи с этим достаточно решить лишь

одно уравнение общего вида.

\[

A_i a_i (\mu_i) \int _0^1 \cos^2 (r {\pi} \xi/2) d\xi = \int _0^1 \cos r ({\pi}\xi/2) d\xi.

\]

Отсюда находим

\begin{equation}

\label{larg10:eq35}

A_i = \pm \frac{4}{r \pi a_i (\mu_i)}, \quad i=1, 2, \dots, 6,

\end{equation}

где знак <<плюс>> "--- для $r=4i-3$; знак <<минус>> "--- для $r=4i-1$.



Отметим, что по формуле \eqref{larg10:eq35} при известных коэффициентах $a_i (\mu_i)$ коэффициенты

$A_i$ могут быть найдены для любого числа приближений $i \in \mathbb N$.



Результаты расчётов безразмерных температур по формуле (\ref{larg10:eq28}) для шести

приближений ($n=6$) в сравнении с решениями, полученными по методам \cite{larg10:3} и

\cite{larg10:4}, представлены на графиках рис.~1,~2. Их анализ позволяет сделать вывод,

что в диапазоне чисел Фурье $0,01 \le {\tt Fo} \le 0,135$ расхождение результатов

составляет не более чем 3{\%}.





\begin{figure}[h!]\centering\footnotesize

\parbox[h]{6.5cm}{\centering \includegraphics[width=6cm]{./00Pic/largina10_1.eps}

\caption{Графики изменения температуры при $\nu =1,0$: 1 "--- по формуле (\ref{larg10:eq28}) (шестое приближение); 2 "--- по методу \cite{larg10:3}}}\hfill \

\parbox[h]{6.5cm}{\centering \includegraphics[width=6cm]{./00Pic/largina10_2.eps}

\caption{Графики изменения температуры при $\nu =1,0$: 1 "--- по формуле (\ref{larg10:eq28}) (шестое приближение); 2 "--- по методу \cite{larg10:4}

(шестое приближение)}}

\end{figure}



Таким образом, на основе использования дополнительных граничных условий

получено аналитическое решение задачи теплопроводности с переменным по

пространственной координате коэффициентом теплопроводности ($\lambda $ "---

экспоненциальная функция координаты). Имеется возможность получения

аналитических решений практически с заданной степенью точности. Решение

имеет простой и удобный для инженерных приложений вид.



Разработанная методика позволяет получать аналитические решения при любой

зависимости теплофизических свойств от пространственной координаты, в том

числе и при одновременном изменении не только коэффициента теплопроводности

от пространственной координаты, но и других теплофизических коэффициентов, а

также даёт возможность получать решения нелинейных задач теплопроводности.



\begin{thebibliography}{9}



\RBibitem{larg10:1}

\by Цой~П.\,В.

\book Методы расчёта задач тепломассопереноса

\publaddr М.

\publ Энергоатомиздат

\yr 1984

\finalinfo 414~с



\RBibitem{larg10:2}

\by Кудинов~В.\,А., Карташов~Э.\,М., Калашников~В.\,В.

\book Аналитические решения задач тепломассопереноса и термоупругости для многослойных конструкций

\publaddr М.

\publ Высш. шк.

\yr 2005

\finalinfo 430~с



\RBibitem{larg10:3}

\by Кудинов~В.\,А., Аверин~Б.\,В., Стефанюк~Е.\,В.

\book Теплопроводность и термоупругость в многослойных конструкциях

\publaddr М.

\publ Высш. шк.

\yr 2008

\finalinfo 391~с



\RBibitem{larg10:4}

\by Ларгина~Е.\,В., Кудинов~И.\,В., Котов~В.\,В., Биктагирова~Е.\,Ю.

\paper Дополнительные граничные условия в задачах теплопроводности с переменными физическими свойствами среды

\jour Вестн. Сам. гос. техн. ун-та. Сер. Математическая

\yr 2009

\issue 2(10)

\pages 48--55



\end{thebibliography}





\makeengtitlem





%\end{document}

